\chapter{Syntax: formulas as objects}
\label{chap:syntax-contexts}

\section{Formulas as typed objects: \texttt{wff}}

A \emph{well-formed formula} in this system is not a bare S-expression
but a typed object that records the theory in which it was formed.  In the REPL it prints as:
\begin{VNBsnip}
#[wff N "formula string"]
\end{VNBsnip}
where \texttt{N} is an internal identifier and the quoted string is the
pretty-printed formula.  Every formula in the system is a \texttt{wff},
and its theory travels with it invisibly.

\texttt{(make-wff expr)} constructs such an object in the current
theory.  \texttt{(free-vars p)} returns the variables that are free in
it.

Two procedures expose the internals of a \texttt{wff}:
\begin{VNBsnip}
(wff->sexp p)        ; => raw S-expression (no output)
(display-contents p) ; prints kind, theory, formula
\end{VNBsnip}
\texttt{wff->sexp} is defined in \texttt{wff.scm};
\texttt{display-contents} is defined in \texttt{contexts.scm}.
Both are available in interactive sessions.

Example (the base theory is \texttt{VNB-SET-THEORY}):
\begin{VNBsnip}
(define p (make-wff "forall([x in A], 0 <= funny-name)"))
; => #[wff 12 "forall([x in a], 0 <= funny-name)"]

(free-vars p)
; => (a funny-name)   ; x is bound by the forall
\end{VNBsnip}

\section{The syntax context}

The destructuring quantifier operates through a \emph{syntax context}:
a table mapping user-facing names to \texttt{nth($k$, $t_0$)}
projections of the implicit tuple instance.  The syntax context is consulted at two points:

\begin{itemize}
\item \textbf{Parser (input):} names in the syntax context are expanded
  to their \texttt{nth} forms before a formula enters the proof state.
\item \textbf{Printer (output):} \texttt{nth} projections of a known
  implicit instance are collapsed back to the declared names when the
  proof state is displayed.
\end{itemize}

The syntax context must travel with the proof state and with
\texttt{wff} objects.  Without the inverse mapping at display
time, the user would see anonymous \texttt{nth} projections instead of
the declared names, defeating the purpose of the mechanism.

The scope of a destructuring quantifier is lexical --- the quantifier
body --- and its basis is the structure of the formula itself.  A second,
\emph{dynamic} discipline once shared this machinery:
\texttt{declare-local-context} conjured a named instance of a structure
into the session until it was undeclared.  It was removed in July 2026,
having acquired no users outside the test suite; see the note at the head
of \texttt{contexts.scm}.

